\chapter{Quantum States and Hilbert Space}

% The braket package ships \ketbra commented out; define it if the preamble has not.
\providecommand{\ketbra}[2]{\ket{#1}\!\bra{#2}}

This chapter develops the mathematical stage on which all of quantum computing
takes place. The single postulate we need is deceptively simple: \emph{the state
of a quantum system is a unit vector in a complex inner-product space.} From this
one idea we recover Dirac's bra--ket notation, the geometry of overlaps and
orthogonality, the description of a single qubit and the Bloch sphere, and --- via
the tensor product --- the exponentially large state space of many qubits, where
entanglement first appears. Everything that follows in later chapters (gates,
circuits, and algorithms) is built on top of this vocabulary.

\section{Finite-Dimensional Hilbert Spaces}

\begin{definition}[Complex inner-product space]\label{l1:def:inner}
Let $\mathcal{H}$ be a complex vector space. An \emph{inner product} on
$\mathcal{H}$ is a map $\langle\,\cdot\,,\,\cdot\,\rangle:\mathcal{H}\times\mathcal{H}\to\C$
satisfying, for all $u,v,w\in\mathcal{H}$ and $a,b\in\C$:
\begin{enumerate}[label=(\roman*)]
  \item \emph{Conjugate symmetry:} $\langle u,v\rangle=\overline{\langle v,u\rangle}$.
  \item \emph{Linearity in the second argument:}
        $\langle u,\,a v+b w\rangle=a\langle u,v\rangle+b\langle u,w\rangle$.
        (Combined with (i), this makes the inner product \emph{conjugate-linear}
        in the first argument.)
  \item \emph{Positive-definiteness:} $\langle u,u\rangle\ge 0$, with equality
        if and only if $u=0$.
\end{enumerate}
\end{definition}

We adopt the physicists' convention that the inner product is linear in the
\emph{second} slot and conjugate-linear in the first; this is the convention that
makes Dirac notation come out cleanly. Concretely, on $\C^{N}$ with
$x=(x_0,\dots,x_{N-1})^{T}$ and $y=(y_0,\dots,y_{N-1})^{T}$,
\begin{equation}
  \langle x,y\rangle \;=\; x^{\dagger}y \;=\; \sum_{i=0}^{N-1}\overline{x_i}\,y_i .
\end{equation}

\begin{definition}[Norm and unit vector]\label{l1:def:norm}
The \emph{norm} induced by the inner product is
$\norm{u}=\sqrt{\langle u,u\rangle}$. A vector $u$ is a \emph{unit vector}
(or \emph{normalized}) if $\norm{u}=1$. Two vectors $u,v$ are \emph{orthogonal}
if $\langle u,v\rangle=0$.
\end{definition}

\begin{definition}[Finite-dimensional Hilbert space]\label{l1:def:hilbert}
A \emph{finite-dimensional Hilbert space} is a finite-dimensional complex
inner-product space $\mathcal{H}$ (in finite dimension completeness is automatic,
so no further condition is needed). If $\dim\mathcal{H}=N$, then choosing an
orthonormal basis gives an isomorphism $\mathcal{H}\cong\C^{N}$ that preserves the
inner product. For a system of $n$ qubits we will find $N=2^{n}$.
\end{definition}

The state-space postulate can now be stated precisely.

\begin{remark}[State-space postulate]
The state of an isolated quantum system is described by a \emph{unit} vector in a
Hilbert space $\mathcal{H}$. The space is closed under linear combination: if
$u,v\in\mathcal{H}$ and $a,b\in\C$, then $a u+b v\in\mathcal{H}$. This closure is
the \emph{superposition principle} --- any linear combination of legal states is
again a legal (unnormalized) state.
\end{remark}

\subsection{The inner product controls all geometry}

\begin{proposition}[Elementary consequences of the axioms]\label{l1:prop:elem}
For all $u,v\in\mathcal{H}$ and $a\in\C$:
\begin{enumerate}[label=(\alph*)]
  \item $\langle u,0\rangle=\langle 0,u\rangle=0$;
  \item $\norm{a u}=\abs{a}\,\norm{u}$;
  \item \emph{(Pythagoras)} if $\langle u,v\rangle=0$ then
        $\norm{u+v}^{2}=\norm{u}^{2}+\norm{v}^{2}$.
\end{enumerate}
\end{proposition}

\begin{proof}
(a) By linearity in the second argument, $\langle u,0\rangle=\langle u,0\cdot 0\rangle=0\,\langle u,0\rangle=0$;
the other equality follows from conjugate symmetry.
(b) $\norm{a u}^{2}=\langle a u,a u\rangle=\overline{a}\,a\,\langle u,u\rangle=\abs{a}^{2}\norm{u}^{2}$,
and take square roots.
(c) Expand $\norm{u+v}^{2}=\langle u+v,u+v\rangle=\norm{u}^{2}+\langle u,v\rangle+\langle v,u\rangle+\norm{v}^{2}$;
the two cross terms vanish because $\langle u,v\rangle=0$ forces
$\langle v,u\rangle=\overline{\langle u,v\rangle}=0$.
\end{proof}

The single most useful inequality in the whole subject is the following.

\begin{theorem}[Cauchy--Schwarz inequality]\label{l1:thm:cs}
For all $u,v\in\mathcal{H}$,
\begin{equation}
  \abs{\langle u,v\rangle}\;\le\;\norm{u}\,\norm{v},
\end{equation}
with equality if and only if $u$ and $v$ are linearly dependent.
\end{theorem}

\begin{proof}
If $v=0$ both sides are $0$ and $u,v$ are trivially dependent, so assume $v\neq0$.
For any $\lambda\in\C$, positive-definiteness gives
\[
  0\;\le\;\norm{u-\lambda v}^{2}
   \;=\;\norm{u}^{2}-\overline{\lambda}\langle v,u\rangle-\lambda\langle u,v\rangle+\abs{\lambda}^{2}\norm{v}^{2}.
\]
Choose $\lambda=\dfrac{\langle v,u\rangle}{\norm{v}^{2}}$ (the value that minimizes the
right-hand side). Then $\lambda\langle u,v\rangle=\dfrac{\langle v,u\rangle\langle u,v\rangle}{\norm{v}^{2}}=\dfrac{\abs{\langle u,v\rangle}^{2}}{\norm{v}^{2}}$, and likewise $\overline{\lambda}\langle v,u\rangle=\dfrac{\abs{\langle u,v\rangle}^{2}}{\norm{v}^{2}}$, so
\[
  0\;\le\;\norm{u}^{2}-\frac{\abs{\langle u,v\rangle}^{2}}{\norm{v}^{2}}.
\]
Rearranging yields $\abs{\langle u,v\rangle}^{2}\le\norm{u}^{2}\norm{v}^{2}$, i.e.
the claim. Equality holds iff $\norm{u-\lambda v}=0$, i.e. $u=\lambda v$, which is
exactly linear dependence.
\end{proof}

\begin{corollary}[Triangle inequality]
$\norm{u+v}\le\norm{u}+\norm{v}$ for all $u,v\in\mathcal{H}$.
\end{corollary}

\begin{proof}
$\norm{u+v}^{2}=\norm{u}^{2}+2\,\mathrm{Re}\,\langle u,v\rangle+\norm{v}^{2}
 \le \norm{u}^{2}+2\abs{\langle u,v\rangle}+\norm{v}^{2}
 \le \norm{u}^{2}+2\norm{u}\norm{v}+\norm{v}^{2}=(\norm{u}+\norm{v})^{2}$,
using $\mathrm{Re}\,z\le\abs{z}$ and Theorem~\ref{l1:thm:cs}. Take square roots.
\end{proof}

\subsection{Geometric meaning of the overlap}

For unit vectors the Cauchy--Schwarz inequality reads $\abs{\langle u,v\rangle}\le1$,
so we may write
\begin{equation}\label{l1:eq:geomdecomp}
  \langle u,v\rangle \;=\; \underbrace{\norm{u}\,\norm{v}}_{\text{length}}\;
  \underbrace{\cos\theta}_{\text{alignment}}\;\underbrace{e^{i\varphi}}_{\text{relative phase}},
  \qquad \cos\theta=\frac{\abs{\langle u,v\rangle}}{\norm{u}\,\norm{v}}\in[0,1].
\end{equation}
The modulus measures how much the two directions coincide ($\cos\theta=1$ means
parallel, $\cos\theta=0$ means orthogonal), while the argument $\varphi$ is a
\emph{relative phase} with no analogue for real vectors. Orthogonality,
$\langle u,v\rangle=0$, is exactly the statement that the two states are perfectly
distinguishable.

\begin{example}[An explicit overlap computation]\label{l1:ex:overlap}
Take the unit vectors in $\C^{2}$
\[
  x=\begin{pmatrix}1\\[2pt]0\end{pmatrix},\qquad
  y=\begin{pmatrix}\tfrac{1}{2}e^{i\pi/6}\\[3pt]\tfrac{\sqrt3}{2}\end{pmatrix}.
\]
One checks $\norm{x}^{2}=1$ and $\norm{y}^{2}=\bigl|\tfrac12\bigr|^{2}+\bigl|\tfrac{\sqrt3}{2}\bigr|^{2}=\tfrac14+\tfrac34=1$.
The overlap is
\[
  x^{\dagger}y=\overline{1}\cdot\tfrac12 e^{i\pi/6}+\overline{0}\cdot\tfrac{\sqrt3}{2}
             =\tfrac12 e^{i\pi/6}.
\]
Comparing with \eqref{l1:eq:geomdecomp}: the length factor is $1$, the alignment is
$\cos\theta=\tfrac12$ (so $\theta=60^{\circ}$), and the relative phase is
$e^{i\pi/6}$ (so $\varphi=30^{\circ}$). The angle and the phase are genuinely
different numbers, which is what makes the two factors independent.
\end{example}

\section{Dirac Notation and Quantum States}

\begin{definition}[Ket and bra]\label{l1:def:braket}
A vector of $\mathcal{H}$ is written as a \emph{ket} $\ket{x}$. To each ket there
corresponds a \emph{bra} $\bra{x}$, the element of the dual space
$\mathcal{H}^{*}$ acting by $\bra{x}:\ket{y}\mapsto\langle x,y\rangle$. In
coordinates, if $\ket{x}$ is the column vector with entries $x_i$, then $\bra{x}$
is the row vector $x^{\dagger}=(\overline{x_0},\dots,\overline{x_{N-1}})$, the
conjugate transpose.
\end{definition}

\begin{definition}[Inner and outer products]\label{l1:def:inout}
The \emph{inner product} of $\ket{x}$ and $\ket{y}$ is the scalar
$\braket{x|y}=\langle x,y\rangle\in\C$. The \emph{outer product} $\ketbra{x}{y}$ is
the linear operator acting as
\[
  \bigl(\ketbra{x}{y}\bigr)\ket{z}=\ket{x}\,\braket{y|z}
  =\braket{y|z}\,\ket{x}.
\]
In coordinates it is the matrix with entries $\bigl(\ketbra{x}{y}\bigr)_{ij}=x_i\,\overline{y_j}$.
\end{definition}

Fix an orthonormal basis $\{\ket{i}\}_{i=0}^{N-1}$, meaning $\braket{i|j}=\delta_{ij}$.
Any ket expands as $\ket{x}=\sum_i x_i\ket{i}$, and applying $\bra{j}$ picks out a
single coefficient:
\begin{equation}\label{l1:eq:coeff}
  \braket{j|x}=\sum_i x_i\braket{j|i}=\sum_i x_i\,\delta_{ji}=x_j .
\end{equation}
So the $j$-th amplitude is literally the overlap $\braket{j|x}$.

\begin{theorem}[Completeness / resolution of the identity]\label{l1:thm:completeness}
If $\{\ket{i}\}_{i=0}^{N-1}$ is an orthonormal basis of $\mathcal{H}$, then
\begin{equation}
  \sum_{i=0}^{N-1}\ketbra{i}{i}\;=\;\id .
\end{equation}
Consequently every operator $A$ satisfies $A=\sum_{i,j}A_{ij}\ketbra{i}{j}$ with
$A_{ij}=\braket{i|A|j}$.
\end{theorem}

\begin{proof}
Let $P=\sum_i\ketbra{i}{i}$. For an arbitrary $\ket{x}=\sum_j x_j\ket{j}$,
\[
  P\ket{x}=\sum_i\ket{i}\braket{i|x}=\sum_i x_i\ket{i}=\ket{x},
\]
using \eqref{l1:eq:coeff}. Since $P\ket{x}=\ket{x}$ for every $\ket{x}$, we have
$P=\id$. For the second statement, insert the identity on both sides of $A$:
$A=\id\,A\,\id=\sum_{i,j}\ket{i}\braket{i|A|j}\bra{j}=\sum_{i,j}A_{ij}\ketbra{i}{j}$.
\end{proof}

\begin{remark}[Projectors]
For a unit vector $\ket{u}$ the operator $P_u=\ketbra{u}{u}$ is the orthogonal
\emph{projector} onto the line spanned by $\ket{u}$: it is Hermitian
($P_u^{\dagger}=P_u$) and idempotent, since
$P_u^{2}=\ket{u}\braket{u|u}\bra{u}=\ket{u}\bra{u}=P_u$ because $\braket{u|u}=1$.
The completeness relation of Theorem~\ref{l1:thm:completeness} decomposes $\id$ into a
sum of orthogonal rank-one projectors, one per basis vector.
\end{remark}

\subsection{The single qubit}

\begin{definition}[Qubit and computational basis]\label{l1:def:qubit}
A \emph{qubit} is a quantum system whose state space is $\mathcal{H}=\C^{2}$. The
distinguished orthonormal basis
\[
  \ket{0}=\begin{pmatrix}1\\0\end{pmatrix},\qquad
  \ket{1}=\begin{pmatrix}0\\1\end{pmatrix}
\]
is called the \emph{computational basis}. A general pure state is
\begin{equation}\label{l1:eq:qubit}
  \ket{\psi}=a\ket{0}+b\ket{1},\qquad a,b\in\C,\qquad \abs{a}^{2}+\abs{b}^{2}=1 .
\end{equation}
The coefficients $a,b$ are called \emph{amplitudes}; $\abs{a}^{2}$ and $\abs{b}^{2}$
are the probabilities of obtaining $0$ or $1$ in a computational-basis measurement.
\end{definition}

\begin{remark}[Global phase is unobservable]\label{l1:rem:global}
Measurement probabilities depend on the amplitudes only through moduli of the form
$\abs{\braket{i|\psi}}^{2}$. Replacing $\ket{\psi}$ by $e^{i\gamma}\ket{\psi}$ for a
real \emph{global phase} $\gamma$ leaves every such probability unchanged, and more
generally leaves $\braket{\phi|A|\psi}$-type quantities invariant when the same
phase appears on both sides. Hence $e^{i\gamma}\ket{\psi}$ and $\ket{\psi}$ describe
the same physical state: genuine states are points of the \emph{projective space}
$\mathbb{P}(\mathcal{H})$, i.e. rays rather than vectors. \emph{Relative} phases
between the terms of a superposition, by contrast, are physical and drive
interference.
\end{remark}

\begin{proposition}[Bloch parametrization]\label{l1:prop:bloch}
Every single-qubit pure state can be written, up to an unobservable global phase, as
\begin{equation}\label{l1:eq:blochform}
  \ket{\psi}=\cos\tfrac{\theta}{2}\,\ket{0}+e^{i\varphi}\sin\tfrac{\theta}{2}\,\ket{1},
  \qquad \theta\in[0,\pi],\quad \varphi\in[0,2\pi).
\end{equation}
The pair $(\theta,\varphi)$ is the pair of spherical angles of a point on the unit
sphere $S^{2}\subset\R^{3}$, the \emph{Bloch sphere}.
\end{proposition}

\begin{proof}
Start from \eqref{l1:eq:qubit} with $\abs{a}^{2}+\abs{b}^{2}=1$. Write the amplitudes in
polar form $a=r_a e^{i\alpha}$ and $b=r_b e^{i\beta}$ with $r_a,r_b\ge0$. Since
$r_a^{2}+r_b^{2}=1$ we may set $r_a=\cos\tfrac{\theta}{2}$ and
$r_b=\sin\tfrac{\theta}{2}$ for a unique $\theta\in[0,\pi]$ (both cosine and sine of
$\theta/2$ are then nonnegative). Thus
\[
  \ket{\psi}=e^{i\alpha}\cos\tfrac{\theta}{2}\ket{0}+e^{i\beta}\sin\tfrac{\theta}{2}\ket{1}
           =e^{i\alpha}\Bigl(\cos\tfrac{\theta}{2}\ket{0}+e^{i(\beta-\alpha)}\sin\tfrac{\theta}{2}\ket{1}\Bigr).
\]
The overall factor $e^{i\alpha}$ is a global phase, which by
Remark~\ref{l1:rem:global} we may discard. Setting $\varphi=\beta-\alpha\ \mathrm{mod}\ 2\pi$
gives \eqref{l1:eq:blochform}. When $\theta\in\{0,\pi\}$ the state is $\ket{0}$ or
$\ket{1}$ and $\varphi$ is irrelevant (the poles of the sphere), consistent with the
identification of states with points of $S^{2}$.
\end{proof}

\begin{remark}[Reading the Bloch sphere]
In this picture $\ket{0}$ and $\ket{1}$ sit at the north and south poles
($\theta=0,\pi$), and \emph{orthogonal} qubit states are \emph{antipodal} points,
not perpendicular ones --- a state and its orthogonal complement are separated by
angle $\theta=\pi$ on the sphere even though they are separated by $90^{\circ}$ in
$\C^{2}$. The equatorial states $\tfrac{1}{\sqrt2}(\ket{0}+e^{i\varphi}\ket{1})$ have
$\theta=\pi/2$ and differ only in their relative phase $\varphi$.
\end{remark}

\begin{example}[Amplitudes and normalization]
The state $\ket{+}=\tfrac{1}{\sqrt2}(\ket{0}+\ket{1})$ has
$\abs{a}^{2}=\abs{b}^{2}=\tfrac12$, so a computational-basis measurement yields $0$ or
$1$ with equal probability $\tfrac12$; its Bloch angles are
$\theta=\pi/2,\varphi=0$. The state $\ket{-}=\tfrac{1}{\sqrt2}(\ket{0}-\ket{1})$ has
the \emph{same} measurement statistics in the computational basis but differs by the
relative phase $\varphi=\pi$; indeed $\braket{+|-}=\tfrac12(1-1)=0$, so $\ket{+}$ and
$\ket{-}$ are orthogonal and perfectly distinguishable by a suitable measurement.
\end{example}

\section{Multi-Qubit Systems and Entanglement}

\begin{definition}[Tensor-product state space]\label{l1:def:tensor}
If systems $A$ and $B$ have state spaces $\mathcal{H}_A$ and $\mathcal{H}_B$, the
composite system has state space the \emph{tensor product}
$\mathcal{H}_A\otimes\mathcal{H}_B$. If $\{\ket{i}_A\}$ and $\{\ket{j}_B\}$ are
orthonormal bases, then $\{\ket{i}_A\otimes\ket{j}_B\}$ is an orthonormal basis of the
product, with inner product
\[
  \bigl(\bra{i}\otimes\bra{j}\bigr)\bigl(\ket{k}\otimes\ket{l}\bigr)
    =\braket{i|k}\,\braket{j|l}.
\]
Hence $\dim(\mathcal{H}_A\otimes\mathcal{H}_B)=(\dim\mathcal{H}_A)(\dim\mathcal{H}_B)$.
For $n$ qubits,
\[
  \mathcal{H}=(\C^{2})^{\otimes n}\cong\C^{2^{n}},\qquad \dim\mathcal{H}=2^{n}.
\]
\end{definition}

We abbreviate $\ket{i}\otimes\ket{j}$ as $\ket{i}\ket{j}$ or $\ket{ij}$. The
computational basis of $n$ qubits is $\{\ket{x}:x\in\{0,1\}^{n}\}$, indexed by the
$2^{n}$ bit strings. The exponential growth of the dimension --- $50$ qubits already
give a space of dimension $2^{50}>10^{15}$ --- is one root of quantum computing's
potential power: a classical description of a general state requires exponentially
many amplitudes.

\begin{definition}[Product and entangled states]\label{l1:def:product}
A state $\ket{\psi}\in\mathcal{H}_A\otimes\mathcal{H}_B$ is a \emph{product state}
(or \emph{separable}) if it can be written as $\ket{\psi}=\ket{\alpha}_A\otimes\ket{\beta}_B$
for some $\ket{\alpha}_A\in\mathcal{H}_A$, $\ket{\beta}_B\in\mathcal{H}_B$. Otherwise
$\ket{\psi}$ is \emph{entangled}.
\end{definition}

\begin{example}[A generic two-qubit state]
Two qubits live in $\C^{4}$ with basis $\ket{00},\ket{01},\ket{10},\ket{11}$. A
general state is
\[
  \ket{\psi}=\sum_{i,j\in\{0,1\}}x_{ij}\ket{ij},\qquad
  \sum_{i,j}\abs{x_{ij}}^{2}=1 .
\]
For instance $\tfrac12(\ket{00}+\ket{01}+\ket{10}+\ket{11})$ is normalized
($4\cdot\tfrac14=1$) and, as the reader can check, factors as
$\tfrac{1}{\sqrt2}(\ket{0}+\ket{1})\otimes\tfrac{1}{\sqrt2}(\ket{0}+\ket{1})$, so it
is a product state.
\end{example}

\begin{definition}[Bell state]\label{l1:def:bell}
The \emph{Bell state} (or EPR pair) is the two-qubit state
\begin{equation}
  \ket{B^{+}}=\tfrac{1}{\sqrt2}\bigl(\ket{00}+\ket{11}\bigr).
\end{equation}
\end{definition}

\begin{theorem}[The Bell state is entangled]\label{l1:thm:bell}
There is no pair of single-qubit states $\ket{\alpha},\ket{\beta}\in\C^{2}$ with
$\ket{B^{+}}=\ket{\alpha}\otimes\ket{\beta}$.
\end{theorem}

\begin{proof}
Suppose for contradiction that such a factorization exists. Write
$\ket{\alpha}=a_0\ket{0}+a_1\ket{1}$ and $\ket{\beta}=b_0\ket{0}+b_1\ket{1}$. Then
\[
  \ket{\alpha}\otimes\ket{\beta}
    =a_0 b_0\ket{00}+a_0 b_1\ket{01}+a_1 b_0\ket{10}+a_1 b_1\ket{11}.
\]
Matching coefficients with $\ket{B^{+}}=\tfrac{1}{\sqrt2}\ket{00}+0\cdot\ket{01}+0\cdot\ket{10}+\tfrac{1}{\sqrt2}\ket{11}$ requires
\[
  a_0 b_0=\tfrac{1}{\sqrt2},\qquad a_0 b_1=0,\qquad a_1 b_0=0,\qquad a_1 b_1=\tfrac{1}{\sqrt2}.
\]
From $a_0 b_0=\tfrac{1}{\sqrt2}\neq0$ we get $a_0\neq0$ and $b_0\neq0$. Then
$a_0 b_1=0$ forces $b_1=0$, while $a_1 b_0=0$ forces $a_1=0$; hence
$a_1 b_1=0\neq\tfrac{1}{\sqrt2}$, a contradiction. Therefore no factorization exists
and $\ket{B^{+}}$ is entangled.
\end{proof}

\begin{remark}[Physical meaning of entanglement]
The correlations of $\ket{B^{+}}$ have no separable description. Measuring the first
qubit in the computational basis yields $0$ or $1$ with probability $\tfrac12$ each;
whatever the outcome, the post-measurement state of the second qubit is forced to the
\emph{same} value ($0\mapsto\ket{0}$, $1\mapsto\ket{1}$). The two outcomes are
perfectly correlated even if the qubits are spatially far apart. This correlation,
inexpressible as any classical mixture of independent single-qubit states, is the
resource behind teleportation, superdense coding, and much of quantum advantage.
\end{remark}

\section{Problems}

\begin{exercise}\label{l1:ex:norm}
Let $\ket{\psi}=2i\ket{0}-\ket{1}\in\C^{2}$ (unnormalized). Find $\norm{\psi}$ and
the normalized state $\ket{\hat\psi}$, and give the probabilities of measuring $0$
and $1$ in the computational basis.
\end{exercise}

\begin{exercise}\label{l1:ex:orth}
Show that $\ket{+}=\tfrac{1}{\sqrt2}(\ket{0}+\ket{1})$ and
$\ket{-}=\tfrac{1}{\sqrt2}(\ket{0}-\ket{1})$ form an orthonormal basis of $\C^{2}$,
and express $\ket{0}$ and $\ket{1}$ in this basis.
\end{exercise}

\begin{exercise}\label{l1:ex:cseq}
Using the equality case of the Cauchy--Schwarz inequality
(Theorem~\ref{l1:thm:cs}), prove that two \emph{unit} vectors $\ket{u},\ket{v}$ satisfy
$\abs{\braket{u|v}}=1$ if and only if $\ket{v}=e^{i\gamma}\ket{u}$ for some real
$\gamma$ --- i.e. they are the same physical state.
\end{exercise}

\begin{exercise}\label{l1:ex:proj}
Let $\ket{u}$ be a unit vector and $P=\ketbra{u}{u}$. Prove that $P^{2}=P$ and
$P^{\dagger}=P$, and compute $\Tr(P)$.
\end{exercise}

\begin{exercise}\label{l1:ex:bloch}
Find the Bloch angles $(\theta,\varphi)$ of the state
$\ket{\psi}=\tfrac{1}{2}\ket{0}+\tfrac{\sqrt3}{2}e^{i\pi/3}\ket{1}$.
\end{exercise}

\begin{exercise}\label{l1:ex:sep}
Decide whether $\ket{\psi}=\tfrac{1}{\sqrt2}(\ket{00}+\ket{01})$ is a product state.
If so, exhibit an explicit factorization; if not, prove it is entangled.
\end{exercise}

\begin{exercise}\label{l1:ex:bellmeas}
For the Bell state $\ket{B^{+}}=\tfrac{1}{\sqrt2}(\ket{00}+\ket{11})$, compute the
probability of obtaining outcome $0$ on the \emph{first} qubit, and give the
resulting (normalized) state of the second qubit.
\end{exercise}

\section*{Solutions}

\begin{solution}[Exercise~\ref{l1:ex:norm}]
$\norm{\psi}^{2}=\abs{2i}^{2}+\abs{-1}^{2}=4+1=5$, so $\norm{\psi}=\sqrt5$. The
normalized state is $\ket{\hat\psi}=\tfrac{1}{\sqrt5}(2i\ket{0}-\ket{1})$. The
probabilities are $\abs{2i/\sqrt5}^{2}=4/5$ for outcome $0$ and
$\abs{-1/\sqrt5}^{2}=1/5$ for outcome $1$; they sum to $1$ as required.
\end{solution}

\begin{solution}[Exercise~\ref{l1:ex:orth}]
$\braket{+|+}=\tfrac12(1+1)=1$ and $\braket{-|-}=\tfrac12(1+1)=1$, so both are unit
vectors. $\braket{+|-}=\tfrac12(\braket{0|0}-\braket{0|1}+\braket{1|0}-\braket{1|1})=\tfrac12(1-0+0-1)=0$,
so they are orthogonal; two orthonormal vectors in a $2$-dimensional space form a
basis. Adding and subtracting, $\ket{0}=\tfrac{1}{\sqrt2}(\ket{+}+\ket{-})$ and
$\ket{1}=\tfrac{1}{\sqrt2}(\ket{+}-\ket{-})$.
\end{solution}

\begin{solution}[Exercise~\ref{l1:ex:cseq}]
For unit vectors, $\norm{u}=\norm{v}=1$, so Cauchy--Schwarz gives
$\abs{\braket{u|v}}\le1$ with equality iff $\ket{u},\ket{v}$ are linearly dependent,
i.e. $\ket{v}=\lambda\ket{u}$ for some $\lambda\in\C$. Taking norms,
$1=\norm{v}=\abs{\lambda}\norm{u}=\abs{\lambda}$, so $\lambda=e^{i\gamma}$ for some
real $\gamma$. Conversely, if $\ket{v}=e^{i\gamma}\ket{u}$ then
$\braket{u|v}=e^{i\gamma}$ and $\abs{\braket{u|v}}=1$. By Remark~\ref{l1:rem:global} the
two are the same physical state.
\end{solution}

\begin{solution}[Exercise~\ref{l1:ex:proj}]
$P^{2}=\ket{u}\braket{u|u}\bra{u}=\ket{u}\bra{u}=P$ since $\braket{u|u}=1$. Taking the
adjoint, $P^{\dagger}=(\ketbra{u}{u})^{\dagger}=\ketbra{u}{u}=P$. For the trace,
using a basis extension or the identity $\Tr(\ketbra{u}{v})=\braket{v|u}$, we get
$\Tr(P)=\braket{u|u}=1$ (a rank-one projector has trace equal to its rank).
\end{solution}

\begin{solution}[Exercise~\ref{l1:ex:bloch}]
Compare with \eqref{l1:eq:blochform}. The $\ket{0}$ amplitude is real and nonnegative,
$\cos\tfrac{\theta}{2}=\tfrac12$, giving $\theta/2=\pi/3$, so $\theta=2\pi/3$. Then
$\sin\tfrac{\theta}{2}=\sin\tfrac{\pi}{3}=\tfrac{\sqrt3}{2}$, matching the modulus of
the $\ket{1}$ amplitude, and its phase is $\varphi=\pi/3$. Hence
$(\theta,\varphi)=(2\pi/3,\ \pi/3)$.
\end{solution}

\begin{solution}[Exercise~\ref{l1:ex:sep}]
Factor out the first qubit:
$\ket{\psi}=\tfrac{1}{\sqrt2}\ket{0}\otimes(\ket{0}+\ket{1})
=\ket{0}\otimes\tfrac{1}{\sqrt2}(\ket{0}+\ket{1})=\ket{0}\otimes\ket{+}$.
This is an explicit product of single-qubit states, so $\ket{\psi}$ is separable
(a product state), not entangled.
\end{solution}

\begin{solution}[Exercise~\ref{l1:ex:bellmeas}]
Write $\ket{B^{+}}=\tfrac{1}{\sqrt2}\ket{0}\ket{0}+\tfrac{1}{\sqrt2}\ket{1}\ket{1}$.
The projector for outcome $0$ on the first qubit is $\ketbra{0}{0}\otimes\id$. The
outcome probability is the squared norm of the (unnormalized) post-measurement vector
$(\ketbra{0}{0}\otimes\id)\ket{B^{+}}=\tfrac{1}{\sqrt2}\ket{0}\ket{0}$, namely
$\bigl\|\tfrac{1}{\sqrt2}\ket{00}\bigr\|^{2}=\tfrac12$. After normalization the state
is $\ket{0}\otimes\ket{0}$, so the second qubit is left in $\ket{0}$ --- perfectly
correlated with the first outcome, as expected for a Bell pair.
\end{solution}
